\chapter{Drift - descrizione a singola particella}

Quando si è andati a studiare come confinare un plasma la prima descrizione a cui si è guardato è stato il comportamento della singola particella del plasma a seconda dei campi magnetici utilizzati a confinarlo.

%%%

Consideriamo $\vec E\perp\vec B$, dall'equazione di Lorentz sappiamo che $m\,\vec a = q\,(\vec E + \frac{1}{c}\vec v \times \vec B)$ .

\[ \vec u\ =\ \vec v - \left(c\vec E\times\vec B/B^2\right)\, ,\qquad \vec{u} '\ =\ \vec{v} '\]

\[m \vec{u} '\ =\ q\vec E\, +\, q \frac{1}{c} \vec u\times \vec B\, +\, q \frac{1}{B^2}\left(\vec E\times \vec B\right)\times \vec B \] 

\[ \left(\vec E\times \vec B\right)\times \vec B\ =\ -\vec B\times \left(\vec E\times \vec B\right)\ =\ -B^2\vec E\, +\, \left(\vec E\cdot \vec B\right)\, \vec B \]

\[ \Rightarrow\qquad m\vec{u}^\circ\ =\ q\,\left[\cancel{\vec E} + \frac{\vec u\times \vec B}{c} + \frac{(\vec E\cdot \vec B)\, \vec B}{B^2} - \cancel{\vec E} \right] \]

\[ \Rightarrow \qquad m\, u_{\parallel} '\ =\ q\, \vec E_{\parallel} \]

\[ \Rightarrow \qquad m \vec u_{\perp}'\ =\ q \frac{\vec u_{\perp} \times \vec B}{c} \]

Osservando il girocentro, la particella non solo prosegue parallelamente al campo B in una spirale, ma ha un drift:

\[ \Rightarrow \qquad \vec v_{gc}\ =\ v_\parallel \hat b\, +\, \frac{c}{B^2} (\vec E\times \vec B) \]

E se ci fosse anche un campo gravitazionale?

\[ \Rightarrow\qquad \vec v_{gc}\ =\ \frac{m}{q}\, \left(\frac{\vec g\times \vec B}{B^2}\right) \]

---

\section{drift gradB? idk}

Poniamo il campo magnetico $\bar B = B_0 (y) \hat e_z$.

Prendiamo L scala caratteristica del campo magnetico ($L\equiv B/\nabla B$), scegliamo un raggio di Larmor $\rho_L \ll L$ e proseguiamo in uno sviluppo nel quale $\epsilon = \rho_L / L$ ($v=v_0 + \bar v_1 \epsilon + \dots$):

\[ \bar B\ =\ B_0 \hat e_z\, +\, (y-y_{g.c.}) \frac{d B_0}{dy}_{|g.c.}\, \hat e_z \]

\[\begin{cases}
	m \dot v_x\ =\ \frac{q}{c}\, v_y\,\left[B\, +\, (y-y_{g.c.}) \frac{d B_0}{dy}_{|g.c.}\right]\\
	m \dot v_y\ =\ -\frac{q}{c}\, v_x\,\left[B\, +\, (y-y_{g.c.}) \frac{d B_0}{dy}_{|g.c.}\right]
\end{cases}\]

\[\text{ordine zero:} \rightarrow \qquad 
\begin{cases}
	m \dot v_{0,x}\ =\ \frac{q}{c}\, B_0\, v_{0,y}\\
	m \dot v_{0,y}\ =\ \frac{q}{c}\, B_0\, v_{0,x}
\end{cases}\]

\[\text{ordine uno:} \rightarrow \qquad 
\begin{cases}
	m \dot v_{1,x}\ =\ \frac{q}{c}\, B_0\, v_{1,y}+\frac{q}{c}\, v_{0,y}(y-y_{g.c.}) \frac{d B_0}{dy}_{|g.c.}\\
	m \dot v_{1,y}\ =\ - \frac{q}{c}\, B_0\, v_{1,x} - \frac{q}{c}\, v_{0,x}(y-y_{g.c.}) \frac{d B_0}{dy}_{|g.c.}
\end{cases}\]

Per vedere l'effetto del drift sul movimento della singola particella dobbiamo fare una media temporale su un periodo molto maggiore del periodo di ciclotrone ($\omega \ll \omega_c$).

...

\[<v_{0,x}(y-y_{gc})>\ =\ \pm \frac{v^2_{\perp}}{2\omega_c}\]

\[\Rightarrow  \qquad 
\bar v_{grad}\ =\ \bar v_{\nabla B}\ =\ \pm \frac{v^2_{\perp}}{2\omega_c} (\frac{B\times \nabla B}{B^2})\]






















 